\chapter{Gestational Diabetes Mellitus}
\index{Gestational Diabetes Mellitus}
\label{sec:Gestational Diabetes Mellitus}

\section{Overview}
Gestational diabetes mellitus (GDM) is glucose intolerance first recognized during pregnancy, affecting 2--10\% of pregnancies. This metabolic disorder involves dysregulation of normal bodily processes, often requiring ongoing monitoring and management. It is increasingly common due to lifestyle and dietary factors. This pregnancy-related condition requires careful monitoring to ensure the safety of both mother and baby. Early detection and management are essential. Metabolic disorders involve disruption of normal biochemical processes essential for health. These conditions may result from enzyme deficiencies, hormonal imbalances, or lifestyle factors. Many metabolic disorders are chronic and require ongoing management. Lifestyle modifications are often the cornerstone of treatment.
\section{Causes}
This condition happens because of insufficient insulin secretion to overcome pregnancy-related insulin resistance mediated by placental hormones (human placental lactogen, cortisol, progesterone). Metabolic disorders arise from dysregulation of normal biochemical pathways. This may result from enzyme deficiencies, hormonal imbalances, or lifestyle factors that overwhelm normal homeostatic mechanisms. The underlying mechanisms involve complex interactions between genetic predisposition and environmental factors. Disruption of normal physiological processes leads to the characteristic features of the condition. Multiple contributing factors may act synergistically to trigger or worsen the condition.
\section{Risk Factors}


\begin{itemize}[noitemsep]
  \item Family history
  \item Obesity and sedentary lifestyle
  \item Poor dietary habits
  \item Advancing age
  \item Hormonal imbalances
  \item Certain medications
  \item Genetic predisposition and family history
  \item Lifestyle factors (diet, exercise, smoking, alcohol)
  \item Environmental exposures
  \item Underlying medical conditions
  \item Medication use
\end{itemize}
\section{Signs and Symptoms}

\subsection{Common symptoms}
\begin{itemize}[noitemsep]
  \item Clinical features are typically asymptomatic
  \item Screening is performed at 24--28 weeks gestation using the 75-g oral glucose tolerance test.
  \item Pain or discomfort in the affected area
  \end{itemize}

\subsection{Emergency warning signs}
\begin{itemize}[noitemsep]
  \item Severe or worsening symptoms of gestational diabetes mellitus require immediate medical evaluation
\end{itemize}
\section{Progression and Stages}
The condition may progress through different stages of severity over time. Early stages often involve mild or subtle symptoms that may be overlooked. Moderate stages involve more noticeable symptoms affecting daily function. Advanced stages may involve significant impairment and complications.
\section{Complications}
\begin{itemize}[noitemsep]
  \item Complications for the mother include preeclampsia and increased cesarean delivery rate
  \item Fetal/neonatal complications include macrosomia, shoulder dystocia, neonatal hypoglycemia, and long-term risk of obesity and T2DM.
  \item Disease progression leading to worsening symptoms
  \item Secondary infections or injuries
  \item Side effects of treatment
  \item Reduced quality of life
  \item Emotional and psychological impact
\end{itemize}
\section{Diagnosis}
Doctors diagnose this based on the oral glucose tolerance test.

\begin{itemize}[noitemsep]
  \item Comprehensive medical history and physical examination
  \item Laboratory tests to assess organ function
  \item Imaging studies to visualize affected structures
  \item Specialized testing depending on the affected system
  \item Consultation with relevant specialists
\end{itemize}
\section{Prevention}
\begin{itemize}[noitemsep]
  \item Maintain a healthy lifestyle with balanced diet and exercise
  \item Avoid known risk factors
  \item Attend regular medical check-ups for early detection
  \item Follow recommended screening guidelines
\end{itemize}
\section{Treatment Options}
\begin{itemize}[noitemsep]
  \item Treatment options include medical nutrition therapy, exercise, and insulin therapy when glycemic targets are not met with lifestyle measures
  \item Metformin may be used in select cases.
  \item Medications to manage symptoms and address underlying mechanisms
  \item Lifestyle modifications and self-care measures
  \item Physical therapy and rehabilitation
  \item Surgical intervention when indicated
  \item Regular monitoring and follow-up care
\end{itemize}
\section{Living with It}
\begin{itemize}[noitemsep]
  \item Follow your treatment plan as prescribed
  \item Maintain regular follow-up with your healthcare provider
  \item Make necessary lifestyle adjustments
  \item Seek support from family, friends, or support groups
  \item Stay informed about your condition
\end{itemize}
\section{Outlook}
\begin{itemize}[noitemsep]
  \item In most cases, the outlook is favorable
  \item GDM usually resolves postpartum, but women with GDM have a 50\% lifetime risk of developing T2DM and should undergo periodic screening.
\end{itemize}
